\documentclass[11pt]{amsart}
\usepackage{geometry}                % See geometry.pdf to learn the layout options. There are lots.
\geometry{letterpaper}                   % ... or a4paper or a5paper or ... 
%\geometry{landscape}                % Activate for for rotated page geometry
%\usepackage[parfill]{parskip}    % Activate to begin paragraphs with an empty line rather than an indent
\usepackage{graphicx}
\usepackage{subfigure}
\usepackage{amssymb}
\usepackage{epstopdf}
\DeclareGraphicsRule{.tif}{png}{.png}{`convert #1 `dirname #1`/`basename #1 .tif`.png}

\newenvironment{packed_enum}{
\begin{enumerate}
  \setlength{\itemsep}{1pt}
  \setlength{\parskip}{0pt}
  \setlength{\parsep}{0pt}
}{\end{enumerate}}

\title{Pancakes and Ternary Strings}
\author{Math 300}
%\date{}                                           % Activate to display a given date or no date

\begin{document}
\maketitle
%\section{x}
%\subsection{}


{\bf  Problem 1}:  Lazy Caterer Problem:  Let $c_n$ denote the maximum number of pieces a circular pancake can be cut into with $n$ cuts.  What are $c_0$, $c_1$, $c_2$,$c_3$, $c_4$?  By looking at the differences $c_n-c_{n-1}$, make a hypothesis for  a recurrence relation for $c_n$. \vfill

{\bf Problem 2}: How many binary strings of length $n$ that contain 00 are? Call this $b_n$. 
\begin{packed_enum}
\item Consider three distinct ways a length $n$ binary string could start:  1,  01, or  00.  For each of theses case, how many ways are there to fill in the rest of the $n-1$ or $n-2$ positions in the string?  
\item One of these three cases is not like the others.  Which one?
\item Use the above answers to write $b_n$ as a second order linear non-homogenous recurrence relation?  What are the initial conditions?
\end{packed_enum}
\vfill

\end{document}  